problem:  
Let \((X,d,\mathfrak{m})\) be a compact, essentially non‑branching metric measure space satisfying the \(CD(K,N)\) condition for some \(K\in\mathbb{R}\), \(1<N<\infty\).  
Assume the following quantitative Euclideanity and quadratic‑defect properties:

1. (**Reifenberg‑type asymptotic Euclideanity**)  
   There exists a function \(\varepsilon:(0,1)\to[0,\infty)\) with \(\varepsilon(r)\to 0\) as \(r\to0\), and an integer \(k\in\{1,\dots,N\}\), such that for every \(x\in X\) and all sufficiently small \(r>0\),
   \[
   d_{mGH}\big((B_r(x),r^{-1}d,\mathfrak{m}_x^r), (B_1(0^k),|\cdot|,\lambda^k)\big)\le \varepsilon(r),
   \]
   where \(\mathfrak{m}_x^r\) is the normalized restriction of \(\mathfrak{m}\) to \(B_r(x)\).

2. (**Approximate parallelogram law for minimal weak upper gradients**)  
   There exists a modulus \(\delta:\mathbb{R}_+\to[0,\infty)\) with \(\delta(s)\to0\) as \(s\to0\) such that for all 1‑Lipschitz functions \(f,g:X\to\mathbb{R}\), the global defect in the quadraticity of the Cheeger energy satisfies  
   \[
   \left|\mathrm{Ch}(f+g)+\mathrm{Ch}(f-g)-2\,\mathrm{Ch}(f)-2\,\mathrm{Ch}(g)\right|\le \delta(r)
   \]
   whenever both \(f\) and \(g\) are supported in a ball \(B_r(x)\) for some \(x\in X\).

**Question:**  
Does the conjunction of these two quantitative conditions necessarily imply that \((X,d,\mathfrak{m})\) is *infinitesimally Hilbertian*—that is, that the Cheeger energy is exactly quadratic, hence \((X,d,\mathfrak{m})\) is an \(RCD(K,N)\) space?  
Equivalently, can there exist a compact \(CD(K,N)\) space that is asymptotically Euclidean in the measured Gromov–Hausdorff sense and whose energy functional satisfies an arbitrarily small parallelogram defect, yet fails the exact \(RCD\) condition?

Why is it a "good" problem:  
This formulation targets the **quantitative stability and rigidity** of the \(RCD(K,N)\) property—the precise threshold where “almost Euclidean” geometry and “almost quadratic” analysis jointly force fully Riemannian structure.  
Current theory guarantees that \(\varepsilon(r)\to0\) (i.e., Euclidean tangents) and exact quadraticity of the Cheeger energy yield \(RCD(K,N)\); the question here asks whether *approximate* versions already suffice.  

A positive solution would establish a **stability principle** asserting that infinitesimal Hilbertianity is robust under small geometric and analytic defects, linking synthetic Ricci bounds to quantitative Reifenberg‑type geometry and convexity of entropy along Wasserstein geodesics.  
A counterexample would reveal a new class of “almost‑Riemannian yet non‑Hilbertian” \(CD\) spaces, exposing subtle interactions between measured Gromov–Hausdorff convergence, Dirichlet form theory, and calculus on nonsmooth spaces.  

The problem is exceptionally simple to state—quantitative Euclideanity + near‑quadraticity → Riemannian?—yet resolving it would likely demand breakthroughs uniting analysis, metric geometry, and stability theory. This balance of clarity, breadth, and foundational depth exemplifies the “narrow entrance, profound depth” character of a good research problem.